\subsection{BlockReduce / Quai Network}

BlockReduce is an architecture using merged mining to scale blockchains through a hierarchical relationship between chains at different levels.

\bquote{
    We argue that the blockchain trilemma is only a true limit in the context were two functions of a blockchain are conflated, specifically consensus and consistency. [\dots]

    [\dots]

    [\dots] how do we disentangle consistency with consensus while maintaining a secure, scalable, decentralized system.

    [\dots]

    Quai Network achieves a scalable PoW architecture because it makes tradeoffs in the CAP theorem trilemma which allows for greater throughput by delaying global consensus. There is an inherent trade-off that is being made. Quai Network can dramatically increase throughput although somewhere between 3–4 orders of increased TPS it will eventually be constrained by network bandwidth, latency, and delay of cross-chain state propagation.

    [\dots]

    Quai Network is an open-source Proof of Work blockchain network utilizing the capabilities of merged mining to increase throughput and security. Users of Quai Network will enjoy fast transaction times without compromising decentralization and security. [\dots]
}{\href{https://perma.cc/CMG8-EPQK?view-mode=server-side}{\emph{Breaking The Blockchain Trilemma} / blog.quai.network} \href{https://archive.ph/nHVsR}{[m]}}


\quote{
    The chain at the top, Prime, has the largest mining reward but also the greatest difficulty and slowest blocks. The next levels, Regions, and zones have sequentially lower difficulty, and faster blocks enabling greater liveliness.
}

\begin{itemize}
    \item due to HLRC, any prime block can out compete \emph{every} block of a higher level (order).
    so a malicious miner can create a prime block that makes stale all region and zone blocks since the last prime block.
    This is why fork rules should never include \texttt{AND} conditions which depend on future blocks -- i.e., a way for a future block to invalidate a past block.
    You might still end up with a consistent outcome, but it means that the work done on blocks of higher level (order) is practically useless.
    \item due to HLRC, miners of prime blocks (i.e., all miners) must only build on valid histories -- \emph{at all levels}.
    Thus, miners must validate $O(n)$ chains.
    \item even if miners didn't have to validate every block, they still need to download every block.
    otherwise a malicious miner can mine a valid prime block that either: has an invalid zone block, or an unavailable region/zone block.
    So downloading all blocks is a minimum.
\end{itemize}
